\paragraph{Carath\'eodory--Schur Lift and Contraction Realization of Scan Spectral Fingerprints: Right Half-Plane Positivity, Minimal Realization, and Local Discriminant Chains}
Continuing with the one-sided spectral fingerprint representation from Proposition \ref{con:xi-2kappa-minimal-complete-samples}:
there exist \(\kappa\) parameters \(r_j\in\mathbb{D}\) and weights \(c_j>0\) such that
\begin{equation}\label{eq:xi-scan-fingerprint-one-sided-exponential-sum}
\boxed{\;
u_n=\sum_{j=1}^{\kappa}c_j\,r_j^n\qquad(n\ge 0)\;}
\end{equation}
(in general position this representation is taken to be minimal: \(\{r_j\}\) are pairwise distinct).

\begin{proposition}[Carath\'eodory Lift of Scan Spectral Fingerprints: Right Half-Plane Positivity and Coefficient Invertibility]\label{prop:xi-scan-fingerprint-caratheodory-lift}
Under the setting of \eqref{eq:xi-scan-fingerprint-one-sided-exponential-sum}, define the analytic function
\begin{equation}\label{eq:xi-scan-fingerprint-carath}
\mathscr{C}_u(z):=\sum_{j=1}^{\kappa}c_j\,\frac{1+r_j z}{1-r_j z},
\qquad |z|<1.
\end{equation}
Then for all \(|z|<1\), the strict positivity \(\Re \mathscr{C}_u(z)>0\) holds.
Moreover, the Taylor coefficients of \(\mathscr{C}_u\) are in one-to-one correspondence with \((u_n)\):
\begin{equation}\label{eq:xi-scan-fingerprint-carath-series}
\boxed{\;
\mathscr{C}_u(z)
=u_0+2\sum_{n\ge 1}u_n\,z^n,\qquad u_0=\sum_{j=1}^{\kappa}c_j>0.\;}
\end{equation}
Therefore, finite-atom scan spectral fingerprints are strictly equivalent to ``finite-pole Carath\'eodory functions'': the sequence \((u_n)_{n\ge 0}\) is uniquely sealed within the protocol-visible domain into a right half-plane positive trajectory \(\mathscr{C}_u\).
\end{proposition}
\begin{proof}
Let \(w=r_j z\). Since \(|r_j|<1\) and \(|z|<1\), we have \(|w|<1\).
The Cayley transform satisfies
\[
\Re\frac{1+w}{1-w}=\frac{1-|w|^2}{|1-w|^2}>0,
\]
so \(\Re\mathscr{C}_u(z)=\sum_j c_j\,\Re\frac{1+r_j z}{1-r_j z}>0\).
Expanding
\(
\frac{1+w}{1-w}=1+2\sum_{n\ge 1}w^n
\)
and substituting \eqref{eq:xi-scan-fingerprint-one-sided-exponential-sum} yields \eqref{eq:xi-scan-fingerprint-carath-series}. \qed
\end{proof}

\begin{proposition}[Uniqueness of Minimal Contraction Realization: Spectral Fingerprints as State Moments of One-Dimensional Defect Contractions]\label{prop:xi-scan-fingerprint-minimal-contraction-realization}
Let \(\mathscr{C}_u\) be as in \eqref{eq:xi-scan-fingerprint-carath}, and define the Schur function
\[
s(z):=\frac{\mathscr{C}_u(z)-\mathscr{C}_u(0)}{\mathscr{C}_u(z)+\mathscr{C}_u(0)},
\qquad |z|<1.
\]
Then \(s\) belongs to the Schur class: \(|s(z)|<1\) and \(s(0)=0\).
Moreover, there exists a (complex) Hilbert space \(\mathcal{H}\), a completely non-unitary contraction operator \(T\) on \(\mathcal{H}\), and a unit vector \(e\in\mathcal{H}\) such that for all \(|z|<1\),
\begin{equation}\label{eq:xi-scan-fingerprint-contraction-herglotz}
\boxed{\;
\mathscr{C}_u(z)=\mathscr{C}_u(0)\,\langle e,(I+zT)(I-zT)^{-1}e\rangle\;}
\end{equation}
yielding the spectral moment identity
\begin{equation}\label{eq:xi-scan-fingerprint-contraction-moments}
\boxed{\;
\frac{u_n}{u_0}=\langle e,T^{n}e\rangle\qquad(n\ge 0).\;}
\end{equation}
Under the minimality condition
\(
\overline{\Span}\{T^n e:n\ge 0\}=\mathcal{H}
\),
this realization \((T,e)\) is unique up to unitary equivalence.
In the finite-atom case (\(\mathscr{C}_u\) is rational and the representation is minimal), one may take \(\dim\mathcal{H}=\kappa\), and \(\mathrm{Spec}(T)=\{r_j\}_{j=1}^{\kappa}\) (counted with algebraic multiplicity).
\end{proposition}
\begin{proof}[Proof outline]
From \(\Re\mathscr{C}_u>0\) and \(\mathscr{C}_u(0)=u_0>0\), the Cayley transform yields \(|s(z)|<1\).
\par
For existence, in the finite-atom representation take \(\mathcal{H}=\CC^\kappa\), let \(T:=\mathrm{diag}(r_1,\dots,r_\kappa)\), and let
\(
e:=(\sqrt{c_1/u_0},\dots,\sqrt{c_\kappa/u_0})^\top
\).
Then \(\|e\|=1\) and
\[
u_0\,\langle e,(I+zT)(I-zT)^{-1}e\rangle
=\sum_{j=1}^{\kappa}c_j\,\frac{1+r_j z}{1-r_j z}
=\mathscr{C}_u(z),
\]
establishing \eqref{eq:xi-scan-fingerprint-contraction-herglotz}.
Substituting \((I+zT)(I-zT)^{-1}=I+2\sum_{n\ge 1}z^nT^n\) and comparing with \eqref{eq:xi-scan-fingerprint-carath-series} yields \eqref{eq:xi-scan-fingerprint-contraction-moments}.
Under the minimal representation, the projection of \(e\) onto each eigenvector direction is nonzero, so the cyclic subspace is the full space.
\par
Uniqueness is a standard result for minimal contraction realizations: if two minimal realizations produce the same \(\mathscr{C}_u\) (equivalently the same \((u_n)\)), then the cyclic representations determined by the moment sequence \(\langle e,T^n e\rangle\) are isomorphic in the unitary sense. \qed
\end{proof}

\begin{corollary}[McMillan Degree Criterion: \texorpdfstring{$\kappa$}{kappa} Equals the Rational Degree of the Generating Function and the Hankel Rank]\label{cor:xi-scan-fingerprint-mcmillan-degree-hankel-rank}
Let the generating function be
\[
U(z):=\sum_{n\ge 0}u_n z^n.
\]
Then from \eqref{eq:xi-scan-fingerprint-one-sided-exponential-sum} one obtains the rational representation
\[
\boxed{\;
U(z)=\sum_{j=1}^{\kappa}\frac{c_j}{1-r_j z}\;}
\]
so the minimal denominator degree (McMillan degree) of \(U\) equals \(\kappa\) in the minimal representation (counted with multiplicity in the repeated root case).
Furthermore, for the block Hankel matrix
\(
\mathsf{H}_N:=(u_{p+q})_{0\le p,q\le N}
\),
there exists \(N_0\) such that for all \(N\ge N_0\), \(\mathrm{rank}(\mathsf{H}_N)=\kappa\) (see Proposition \ref{con:xi-scan-hankel-rank-equals-defect}).
Therefore, the defect degree \(\kappa\) simultaneously equals:
(i) the minimal rational realization degree;
(ii) the minimal contraction realization dimension (Proposition \ref{prop:xi-scan-fingerprint-minimal-contraction-realization});
(iii) the stabilized Hankel rank.
\end{corollary}
\begin{proof}[Proof outline]
The rational representation is obtained by termwise summation of the geometric series \(1/(1-rz)=\sum_{n\ge 0}r^n z^n\).
The consistency of the denominator degree with \(\kappa\) in the minimal representation follows from the number of poles (counted with multiplicity).
The Hankel rank conclusion has been established by the minimal realization theory of \(\sum_j c_j r_j^n\) (Theorem \ref{thm:hankel-rank-minimal} and Proposition \ref{con:xi-scan-hankel-rank-equals-defect}). \qed
\end{proof}

\begin{proposition}[Uniqueness of Depth--Height Tensor Factorization: Irreducible Direct Product of Modulus Orbits and Phase Carriers]\label{prop:xi-scan-depth-height-tensor-factorization}
Write \(r_j=\rho_j e^{-\mathrm{i}\theta_j}\) (\(\rho_j:=|r_j|\in(0,1)\)).
Then the spectral fingerprint decomposes as
\[
u_n=\sum_{j=1}^{\kappa}c_j\,\rho_j^{n}\,e^{-\mathrm{i}n\theta_j}.
\]
Define the pure depth modulus orbit
\[
v_n:=\sum_{j=1}^{\kappa}c_j\,\rho_j^{n}\qquad(n\ge 0).
\]
Then \(\{(\rho_j,c_j)\}\) (equivalently any single-variable depth orbit) can only recover the modulus decay structure, while the phase carrier \(\{e^{-\mathrm{i}n\theta_j}\}\) constitutes an irreducible independent input: there exists no protocol-internal channel that can endogenously recover the phase \(\theta\) from the depth-side orbit alone; any construction attempting to eliminate this direct product splitting must degenerate under auditable closure into congruence folding or irrecoverable collapse.
This proposition is a direct specialization of Proposition \ref{prop:xi-depth-height-minimal-coupling} to the discrete address spectral fingerprint regime.
\end{proposition}
\begin{proof}
From \(|r_j|=\rho_j\), the sequence \((v_n)\) equivalently encodes the amplitude spectrum of \(\{(\rho_j,c_j)\}\); the \(\theta_j\) enter only through the characters \(n\mapsto e^{-\mathrm{i}n\theta_j}\).
Proposition \ref{prop:xi-depth-height-minimal-coupling} asserts that the amplitude spectrum and phase spectrum form a minimal direct product coupling at the protocol level, and the height (phase) must be externalized through an independent carrier; restricting the role of \(e^{-\mathrm{i}\gamma_j n\Delta}\) to lattice characters yields the present proposition. \qed
\end{proof}

\begin{proposition}[``Coefficient 4'' in the Spectral Fingerprint Side: Double Cayley Origin of the Factor \texorpdfstring{$2$}{2} Composed Twice]\label{prop:xi-scan-coefficient-four-double-cayley}
In the Carath\'eodory lift of scan spectral fingerprints, the fundamental kernel satisfies
\[
\frac{1+w}{1-w}=1+2\sum_{n\ge 1}w^n\qquad(|w|<1),
\]
so the coefficient prefactor in \(\mathscr{C}_u(z)\) is strictly \(2\) (\eqref{eq:xi-scan-fingerprint-carath-series}).
On the other hand, in the disk/half-plane normalization at the endpoint \(-1\), the conformal factor of the Cayley transform is \(2\), and its square appears as the density pullback constant strictly equal to \(4\) in the boundary Poisson/Busemann coordinates (Proposition \ref{con:xi-four-from-cayley-conformal-factor}).
Therefore, in the composite chain ``spectral fingerprint \(\to\) Carath\'eodory \(\to\) boundary density,'' all occurrences of the constants \(2\)/\(4\) arise from normalization sealing and carry no new dynamical degrees of freedom.
\end{proposition}
\begin{proof}
The first part follows directly from the power series expansion of the Cayley transform.
The second part has been established by Proposition \ref{con:xi-four-from-cayley-conformal-factor}: \(4\) is the square of the conformal factor \(2\), and it appears in auditable invariants only as a uniform rescaling that cancels in ratio coordinates.
Combining the two parts yields the conclusion. \qed
\end{proof}

\paragraph{Pure Depth Subclass: Low-Order Hankel Discriminant Chains and Stepwise Source Count Lower Bounds}
In the following, the phase carriers are collapsed to real-axis moduli: assume \(r_j=\rho_j\in(0,1)\subset\RR\) while keeping \(c_j>0\), so that \((u_n)\) is the power moment sequence of a finite positive measure on \((0,1)\).

\begin{proposition}[Minimal Algebraic Certificate for the Two-Source Threshold: First Discriminant for \texorpdfstring{$\kappa\ge 2$}{kappa>=2}]\label{prop:xi-scan-two-source-min-discriminant}
In the pure depth subclass, define the second-order Hankel discriminant
\[
\Delta_1
:=\det\begin{pmatrix}u_0 & u_1\\[2pt] u_1 & u_2\end{pmatrix}
=u_0u_2-u_1^2.
\]
Then \(\Delta_1\ge 0\) necessarily holds, and there exists an explicit decomposition
\[
\boxed{\;
\Delta_1=\sum_{1\le j<k\le \kappa}c_jc_k(\rho_j-\rho_k)^2\; }.
\]
Hence, in the minimal representation regime (\(\rho_j\) pairwise distinct), the following are equivalent:
\[
\boxed{\;\kappa\ge 2\quad\Longleftrightarrow\quad \Delta_1>0.\;}
\]
\end{proposition}
\begin{proof}
Expanding yields
\[
u_0u_2-u_1^2
=\Bigl(\sum_j c_j\Bigr)\Bigl(\sum_k c_k\rho_k^2\Bigr)-\Bigl(\sum_j c_j\rho_j\Bigr)^2
=\sum_{j,k}c_jc_k(\rho_k^2-\rho_j\rho_k).
\]
Symmetrizing and rearranging into \(\sum_{j<k}c_jc_k(\rho_j-\rho_k)^2\) yields the decomposition, whence \(\Delta_1\ge 0\), and strictly positive when \(\rho_j\) are pairwise distinct and \(\kappa\ge 2\).
Conversely, if \(\Delta_1=0\), the above expression forces all \(\rho_j\) to be equal, so the minimal representation degenerates to a single source. \qed
\end{proof}

\begin{proposition}[Two-Source Persistence and Certificate Stability: Condition Number of the Two-Source Discriminant is Controlled Solely by the Spectral Gap]\label{prop:xi-scan-two-source-stability-gap}
Remaining in the pure depth subclass, order so that \(\rho_1>\rho_2\ge\cdots\ge\rho_\kappa\), and define the spectral gap ratio
\[
\eta:=\frac{\rho_2}{\rho_1}\in(0,1).
\]
Then, in a parameter domain with fixed total mass \(u_0=\sum_j c_j\) and non-degenerate weights (e.g., \(c_1,c_2\) have a uniform positive lower bound),
any perturbation driving the ``two-source'' discriminant \(\Delta_1\) from a positive value toward \(0\) inevitably passes through an ill-conditioned region controlled by \((1-\eta)\);
more quantitatively, the Lipschitz condition number for the local inversion of the two-source separation quantity \(\rho_1-\rho_2=\rho_1(1-\eta)\) diverges at least at the scale \((1-\eta)^{-1}\).
\end{proposition}
\begin{proof}[Proof outline]
By Proposition \ref{prop:xi-scan-two-source-min-discriminant},
\(
\Delta_1\ge c_1c_2(\rho_1-\rho_2)^2
\).
Hence, when \(c_1,c_2\) have a uniform positive lower bound,
\(
\rho_1-\rho_2\le C\sqrt{\Delta_1}
\)
and the inverse derivative scale satisfies
\(
\mathrm{d}(\rho_1-\rho_2)\sim \Delta_1^{-1/2}\mathrm{d}\Delta_1
\);
when \(\rho_2\uparrow\rho_1\) (i.e., \(1-\eta\downarrow 0\)), \(\Delta_1\downarrow 0\) is necessarily triggered, and the condition number diverges at least as \((\rho_1-\rho_2)^{-1}\asymp (1-\eta)^{-1}\). \qed
\end{proof}

\begin{proposition}[From Two Sources to Multiple Sources: Low-Order Discriminant Chain Yields Stepwise Lower Bounds on \texorpdfstring{$\kappa$}{kappa}]\label{prop:xi-scan-hankel-determinant-chain-kappa}
In the pure depth subclass, for \(m\ge 0\) define the Hankel discriminant chain
\[
\Delta_m:=\det(u_{p+q})_{0\le p,q\le m}.
\]
Then \(\Delta_m\ge 0\) for all \(m\).
Moreover, in the minimal representation regime (\(\rho_j\) pairwise distinct), the following stepwise equivalences hold:
\[
\boxed{\;
\kappa\ge m+1\quad\Longleftrightarrow\quad \Delta_m>0,\qquad
\kappa\le m\quad\Longleftrightarrow\quad \Delta_m=0.\;}
\]
Therefore, \(\{\Delta_m\}\) constitutes a completely local ``source count lower bound certificate chain'': without inverting for \(\{\rho_j\}\), one need only verify the strict positivity of finite-order Hankel determinants stepwise to establish within the protocol-visible domain the fact ``at least \(m+1\) sources''; its failure (\(\Delta_m=0\)) corresponds to the structural recovery of the \((m+1)\)-th degree of freedom.
\end{proposition}
\begin{proof}[Proof outline]
Let \(v(\rho):=(1,\rho,\dots,\rho^{m})^\top\in\RR^{m+1}\).
Then the Hankel block can be written as a Gram-type decomposition
\[
(u_{p+q})_{0\le p,q\le m}
=\sum_{j=1}^{\kappa}c_j\,v(\rho_j)v(\rho_j)^\top\succeq 0,
\]
so \(\Delta_m\ge 0\).
If \(\kappa\le m\), the rank of the above matrix does not exceed \(\kappa\), so \(\Delta_m=0\).
If \(\kappa\ge m+1\) and \(\rho_j\) are pairwise distinct, then there exist \(m+1\) nodes whose corresponding Vandermonde vectors are linearly independent, the Gram matrix has full rank, and hence \(\Delta_m>0\).
The equivalence can also be given by the explicit Vandermonde expansion:
\[
\Delta_m=\sum_{1\le j_0<\cdots<j_m\le \kappa}
\Bigl(\prod_{\ell=0}^{m}c_{j_\ell}\Bigr)
\prod_{0\le p<q\le m}(\rho_{j_q}-\rho_{j_p})^2,
\]
where each term is nonneg and at least one term is strictly positive when there exist \(m+1\) distinct nodes. \qed
\end{proof}

\endinput
